\section{What directs verification?}\label{sec:directs}

\paragraph{The association.} Across \ExpOneN{} pre-registered episodes,
allocation tracks the agent's first-pass plan: the corrupted memory is verified
in \AlignedK/\AlignedN{} episodes whose plan it backs and \MisK/\MisN{}
otherwise. A further \SixEpisodes{} episodes over \SixModels{} models from
\SixOrgs{} organizations, in two unrelated domains built isomorphic to the
first, reproduce it: across that whole set the plan-backed memory is verified
in \SixAllOnPathK/\SixAllOnPathN{} episodes whose plan it backs, so the
signature is not a property of one scenario or one laboratory.

The grid also locates where the signature weakens, and the boundaries are as
informative as the replications. At a budget of one, the rate at which any
credit reaches the plan falls from \SixPlanTwoPct\% to \SixPlanOnePct\%
(\SixPlanOneK/\SixPlanOneN{} episodes) --- \emph{always checked} is a property
of having a spare credit, not of the plan. Two of eight dose-by-world cells fall below the registered per-cell floor
(\SixHTwentyFiveLowCells). Pooling every budget-one episode where the plan-backed
and the topically-salient memory differ, the credit goes to the plan in
\SixChoicePlanK/\SixChoicePlanN{} and to salience in
\SixChoiceSalientK/\SixChoiceSalientN{} --- but in one cell
(\SixSalienceWonCell, $n=\SixSalienceWonN$) salience wins outright,
\SixSalienceWonSalPct\% against \SixSalienceWonPlanPct\%. The signature is strong,
replicated, and not universal --- which is what an allocation policy, rather
than a law, should look like.

None of this identifies a cause. The plan and the lookups are named in the same
response, so an agent that simply writes a coherent answer would produce the
same table.

\paragraph{The manipulation.} We therefore assign the plan exogenously and
remove the agent's ability to narrate it. Before the memories are shown, a
provisional direction is set by leadership, randomised between a pricing plan
and an onboarding plan; the agent may confirm or replace it. In the
\emph{verify-only} elicitation the \texttt{intended\_action} field is
\textbf{removed from the response schema}, so the model reports which sources
it wants and nothing else. A third arm assigns no plan at all. \SevenN{}
episodes, \SevenErr{} errors.

\begin{center}\itshape
Assigning a pricing plan moves the single verification credit onto the pricing
memory in \SevenPriK/\SevenPriN{} episodes, against \SevenOnbK/\SevenOnbN{}
under an onboarding plan: \textbf{\SevenDelta{} points}.
\end{center}

\noindent The direction is positive in \SevenPos{} of \SevenModels{} models
(CMH stratified by model, $\chi^2 = \SevenCMH$, $\SevenCMHP$; per-model rates
in Appendix~\ref{app:models}). The mirror contrast runs the same way: an
\emph{onboarding} memory is verified in \SevenMirrorOnbK/\SevenMirrorOnbN{}
episodes under an onboarding plan against \SevenMirrorPriK/\SevenMirrorPriN{}
under a pricing plan, \SevenMirror{} points. The assignment \emph{redirects}
attention rather than merely suppressing one memory.

Against the no-plan baseline (\SevenNoneK/\SevenNoneN) the assignment is
asymmetric: a pricing plan raises target verification by \SevenPriVsNone{}
points and an onboarding plan lowers it by \SevenOnbVsNone{}. Most of the
\SevenDelta{} is the plan steering attention \emph{away}.

\paragraph{What this does and does not rule out.} Removing the action field
does not shrink the effect; verify-only \emph{exceeds} joint elicitation by
\SevenElicitGap{} points. This breaches our registered H32, which required the
two elicitations to agree within 15 points; we report the breach, and note that
its direction is the opposite of what response coherence predicts.

Removing the field defeats the \emph{narration} route --- the agent cannot
write a rationale matching an action it never states. It does not defeat a
residual \textbf{topical priming} account: naming \texttt{promotional\_pricing}
in the assignment makes pricing salient as well as making it the plan, and this
design cannot separate the two. We flag this rather than argue it away; the
cell that would settle it mentions the topic while denying the agent
discretion, and we did not run it.

A descriptive observation cuts both ways and we give both halves. In the joint
elicitation --- the only arm where an action is observable ---
\SevenOverturnRate\% of agents assigned a pricing plan overturn it. Those who
keep the plan verify the pricing memory in \SevenKeptK/\SevenKeptN{} episodes;
those who abandon it still do so in \SevenOverturnK/\SevenOverturnN. Coherence
plainly accounts for much of the joint-arm allocation, and a substantial
fraction survives abandoning the plan. This conditions on a post-treatment
variable and is offered as description, not evidence.

We read the randomised contrast as identifying the upstream half of the chain:
\textbf{an exogenous change to the agent's assigned working plan changes which
inherited beliefs it verifies}. We do not read it as isolating internal plan
state. The intervention bundles working-plan state with the topical salience of
naming the plan and with directive framing --- the direction is set by
leadership --- and the experiment identifies the assignment, not the
independent contribution of those components.
